\section{Purpose and geometric point of view}

Current-plane plots for a \gls{pmsm} commonly show a current circle, a voltage
ellipse, constant-torque curves, and optimal trajectories.  Calling the
voltage boundary an ellipse before inspecting its equation, however, hides the
reason it is closed and the assumptions under which that statement fails.
This paper therefore starts from the voltage map and lets its algebra determine
the conic.  The machine conventions follow standard rotor-oriented models
\cite{binder2017,schroeder2021}; the geometric analysis uses elementary
properties of real symmetric matrices \cite{horn2013}.

The central analogy is a quadratic cost landscape.  A positive-definite
quadratic form is a bowl.  Its level curves are ellipses; its eigenvectors are
directions along which the bowl rises without lateral coupling; and its
eigenvalues measure steepness.  Translation moves the bowl, and rotation
changes its orientation, but neither operation changes its intrinsic conic
class.

\begin{figure}[tb]
\centering
\begin{tikzpicture}[scale=.9]
 \begin{scope}[xshift=-3.3cm]
  \draw[->] (-2.1,0)--(2.2,0) node[right]{$u_d$};
  \draw[->] (0,-2.1)--(0,2.2) node[above]{$u_q$};
  \draw[constraint curve] (0,0) circle (1.55);
  \node[annotation] at (0,-2.35){voltage-space norm};
 \end{scope}
 \draw[-{Latex[length=3mm]},very thick,gray] (-.7,0)--(.7,0)
   node[midway,above]{inverse affine map};
 \begin{scope}[xshift=3.3cm,rotate=18]
  \draw[->,rotate=-18] (-2.3,0)--(2.4,0) node[right]{$i_d$};
  \draw[->,rotate=-18] (0,-2)--(0,2.2) node[above]{$i_q$};
  \draw[constraint curve] (.25,-.15) ellipse (1.9 and 1.05);
  \node[annotation,rotate=-18] at (.2,-2.25){current-space conic};
 \end{scope}
\end{tikzpicture}
\caption{A Euclidean voltage circle pulls back through an invertible affine
machine map to a translated and generally rotated ellipse.}
\label{fig:circle-ellipse}
\end{figure}


The derivation repeatedly separates five questions: what algebraic operation
is performed, why it is valid, what it means geometrically, what physical term
caused it, and what an engineer can infer from it.  The model is linear and
steady state.  Saturation, cross-saturation, iron loss, inverter harmonics, and
transient flux derivatives are outside its scope; their omission is an
explicit modeling decision, not a claim that real machines are exact quadrics.

Accordingly, this is a learning paper rather than a compressed derivation
note.  Important conclusions are approached from several directions.  The
voltage boundary is first expanded component by component, then classified as
a conic, interpreted as an inverse image of a circle, diagonalized through an
eigenbasis, and finally read as the gain geometry of an SVD.  These routes are
kept together because each answers a different question: algebra checks the
coefficients, conic theory names the shape, the inverse-image view explains
why it appears, eigenvectors expose its axes, and singular values connect
those axes to physical voltage gain.
